% =============================================================================
% quiz template — single in-class quiz with optional bank workflow
% =============================================================================
%
% Compile with:
%   lualatex quiz-template.tex
%
% This template demonstrates BOTH styles in one file:
%   1. Inline \question text (the original quiz pattern)
%   2. \getproblem retrieval from a problem bank
%
% Keep whichever you prefer; mix and match if useful.

\documentclass[quiz-number = 1]{quiz}
% Optional: replace the boxed instructions with a separate file by adding
%   quiz-instructions-file = group-quiz-instructions
% to the options above (filename without .tex; \input'd at \maketitle).

% ---- Inline bank entries (or move into a separate `quiz_bank.tex` and
% ---- \loadbank{quiz_bank} below) -------------------------------------
% Defined with the \begin{problem} environment: the stem, then the answer in a
% \begin{solution}...\end{solution} block.  Because they live in this file,
% SyncTeX inverse search lands
% here — double-clicking the typeset problem in the PDF jumps to its
% \begin{problem} block below.  Definitions may go in the preamble (as here) or
% the document body.
\begin{problem}{eval-poly}[topic=eval, diff=easy]
	Evaluate $3 \cdot (-2)^2 + 5 \cdot 4 - 3 \cdot 3^2$.
	\begin{solution} $5$
	\end{solution}
\end{problem}

\begin{problem}{linear-eq}[topic=algebra, diff=easy]
	Solve $2x + 5 = 11$ for $x$.
	\begin{solution} $x = 3$
	\end{solution}
\end{problem}

% Alternatively, keep the bank in its own file:
% \loadbank{quiz_bank.tex}

\begin{document}
\maketitle

\begin{questions}

% --- (a) Bank-backed: pulled from \begin{problem} above ---------------
\question \getproblem{eval-poly}
	\begin{parts}
		\part[3] (Show your arithmetic.)
			\vspace{\stretch{1}}
		\part[2] State the order of operations rules you used.
			\vspace{\stretch{1}}
	\end{parts}

% --- (b) Bank-backed: random match from a key=value query -------------
\question \getproblem{topic=algebra, diff=easy}
	\vspace{\stretch{1}}

% --- (c) Inline (original quiz style — still works) -------------------
\question Solve the inequality $|x - 2| < 5$.
	\vspace{\stretch{1}}

\extracredit[2]{Find a real number $x$ such that $x^x = 1/4$.}
	\vspace{\stretch{1}}

\end{questions}

\end{document}
